\documentclass[11pt]{article}
% Unicode fonts, because the statement and the prose carry Lean's symbols
% (ℝ, ℚ, ∀, ≠, √, ⟨⟩) and the Latin Modern faces under T1 have no glyph for
% them: Tectonic then logs "Missing character" and emits a PDF whose exact
% Lean statement has silently lost its quantifier and its domains. DejaVu
% covers every symbol the recorded runs used. Named by file, not by family:
% a family name is resolved through the host's fontconfig, which served
% /usr/share/fonts on the machine that recorded the third staged run and
% made Tectonic warn that the build may not be reproducible, whereas a file
% name is served by the pinned bundle, whose digest the manifest records.
\usepackage{fontspec}
\setmainfont{DejaVuSans}[Extension=.ttf, UprightFont=*, BoldFont=*-Bold, ItalicFont=*-Oblique, BoldItalicFont=*-BoldOblique]
\setmonofont{DejaVuSansMono}[Extension=.ttf, UprightFont=*, BoldFont=*-Bold, ItalicFont=*-Oblique, BoldItalicFont=*-BoldOblique]
\usepackage[margin=1in]{geometry}
\usepackage{fancyvrb}
\usepackage{xcolor}
\usepackage{enumitem}
\usepackage[hidelinks]{hyperref}

\definecolor{HardyBlue}{RGB}{24,64,112}
\setlength{\parindent}{0pt}
\setlength{\parskip}{0.7em}
\setlist[itemize]{leftmargin=1.5em}

\begin{document}

\begin{center}
{\LARGE\color{HardyBlue}\bfseries The Irrationality of √2 + √3\par}
\vspace{1em}
Formal status: Kernel verified\\
Faithfulness: User approved; independently reviewed by claude-opus-5 on claude (agreed)\\
Informal completeness: Not independently assessed\\
Document status: TeX compiled
\end{center}

\hrule
\vspace{1em}

\section*{Theorem}
The real number √2 + √3 is irrational; that is, there is no rational number q with q = √2 + √3.

Formally (Lean 4 / Mathlib):

    theorem irrational\_sqrt\_two\_add\_sqrt\_three : Irrational (Real.sqrt 2 + Real.sqrt 3)

Here `Irrational x` is Mathlib's definition `x ∉ Set.range ((↑) : ℚ → ℝ)`, i.e. x is not the image of any rational number under the canonical embedding ℚ → ℝ.

\section*{Approved interpretation}
\begin{itemize}
\item The real numbers ℝ, with its standard complete ordered field structure and the nonnegative square root function Real.sqrt : ℝ → ℝ
\item The rational numbers ℚ, embedded into ℝ by the canonical cast Rat.cast : ℚ → ℝ
\item Numeric literals 2 and 3 elaborated as elements of ℝ via OfNat
\item No user-visible quantifiers: the statement is a closed proposition about one specific real number, so the theorem takes no binders at all.
\item The quantification hidden inside 'irrational' is universal-negative: for every q : ℚ, (q : ℝ) ≠ Real.sqrt 2 + Real.sqrt 3, i.e. the number does not lie in Set.range ((↑) : ℚ → ℝ).
\item None; the claim is unconditional and needs no hypotheses beyond Mathlib's standard theory of ℝ, ℚ, and Real.sqrt.
\item Real.sqrt is total in Mathlib (returning 0 on negative inputs); since 2 ≥ 0 and 3 ≥ 0, Real.sqrt 2 and Real.sqrt 3 are the genuine nonnegative square roots.
\item Availability of Mathlib's Irrational predicate (Mathlib.Analysis.SpecialFunctions.Pow / Mathlib.RingTheory.Int.Basic import chain, e.g. via `import Mathlib`).
\item Repair of the previous non-elaborating signature: the binders field previously carried the parenthesized prose '(no binders; closed statement)', which is not valid Lean binder syntax and would be spliced directly into `theorem <name> <binders> : <prop>`, causing elaboration to fail. The binders field is now the empty string, yielding `theorem irrational\_sqrt\_two\_add\_sqrt\_three : Irrational (Real.sqrt 2 + Real.sqrt 3)`.
\item Fully qualified names are used (Real.sqrt, Irrational) so the statement elaborates without relying on `open Real`; no notation such as √ is used, avoiding dependence on a scoped notation being open.
\item 'sqrt(2)' and 'sqrt(3)' are the principal (nonnegative) real square roots, formalized as Real.sqrt 2 and Real.sqrt 3 — not two-valued algebraic roots and not elements of an abstract field extension.
\item 'is irrational' is Mathlib's predicate Irrational x, definitionally x ∉ Set.range ((↑) : ℚ → ℝ), which is exactly 'not equal to the cast of any rational'.
\item '+' is addition in ℝ, and the whole sum is the argument of Irrational, so the claim is about the single number Real.sqrt 2 + Real.sqrt 3 — neither weakened (e.g. to irrationality of Real.sqrt 2 alone) nor strengthened (e.g. to linear independence over ℚ or a general theorem about sums of square roots).
\end{itemize}

\section*{Human-readable proof}
Overview
--------
The sum of two irrational numbers need not be irrational (for instance √2 + (−√2) = 0), so irrationality of √2 and of √3 separately does not settle the question. The argument below instead assumes rationality of the sum and then *solves algebraically for √2*, reducing the statement to the single classical fact that √2 is irrational.

Setup
-----
Suppose, for contradiction, that √2 + √3 is rational. By the definition of `Irrational`, this means there exists q ∈ ℚ whose image in ℝ satisfies

    (q : ℝ) = √2 + √3.                                          (1)

We use two facts about the real square root, both instances of Mathlib's `Real.sq\_sqrt` (valid because the radicands are nonnegative):

    (√2)² = 2,        (√3)² = 3.                                 (2)

Step 1: q ≠ 0
-------------
Since 2 > 0 and 3 > 0, both √2 > 0 and √3 > 0, hence √2 + √3 > 0. By (1), (q : ℝ) > 0, so in particular (q : ℝ) ≠ 0. (In the formal proof this positivity is discharged by the `positivity` tactic, which knows that √x > 0 for x > 0.)

Step 2: isolate and square
--------------------------
Rearranging (1) gives

    √3 = (q : ℝ) − √2.

Squaring both sides and substituting (√3)² = 3 from (2):

    3 = ((q : ℝ) − √2)² = q² − 2q·√2 + (√2)² = q² − 2q·√2 + 2,

where the last equality uses (√2)² = 2. Rearranging the resulting identity yields the key linear relation for √2:

    2q·√2 = q² − 1.                                              (3)

Note that this step is purely algebraic: no sign or branch considerations are needed, because we only expand a square and substitute the two identities in (2).

Step 3: solve for √2 and conclude
---------------------------------
By Step 1 we have q ≠ 0, so 2q ≠ 0 and we may divide (3) by 2q:

    √2 = (q² − 1) / (2q).

The right-hand side is the image in ℝ of the rational number (q² − 1)/(2q) ∈ ℚ — the field operations on ℚ commute with the embedding ℚ → ℝ, so the cast of the rational quotient equals the quotient of the casts. Hence √2 lies in the range of ℚ → ℝ, i.e. √2 is rational.

This contradicts the classical theorem that √2 is irrational (Mathlib's `irrational\_sqrt\_two`, itself a consequence of the fact that 2 is prime and therefore not a perfect square).

Therefore no rational q with (q : ℝ) = √2 + √3 exists, and √2 + √3 is irrational. ∎

Remarks
-------
• A common alternative route squares the sum directly: (√2 + √3)² = 5 + 2√6, so rationality of the sum would force √6 to be rational, contradicting the irrationality of √6. The route taken here is shorter to formalize because it needs only the irrationality of √2, a named lemma in Mathlib, rather than a fresh non-square argument for 6.
• The same technique generalizes: if a and b are distinct nonnegative rationals with √a irrational, the identity above shows √a + √b is irrational, since assuming rationality of the sum lets one solve for √a as a rational expression in the assumed rational value.

Formal status
-------------
The Lean statement above was rebuilt from the Frozen Claim, checked by a fresh Lean kernel using Mathlib, and its axiom dependencies were audited by Hardy's independent FinalVerifier: the formal proof is kernel verified. The formal script follows exactly the three steps above — positivity of the sum to get q ≠ 0, squaring to obtain relation (3) (closed by `nlinarith` from the two square-root identities), then exhibiting (q² − 1)/(2q) as a rational witness for √2 and appealing to `irrational\_sqrt\_two`.

This prose exposition is a human-readable rendering of that formal argument; the exposition itself has not been independently assessed, and the machine-checked guarantee attaches to the Lean statement and proof, not to the wording here.

\section*{Known gaps}
\begin{itemize}
\item None in the formal development: the Lean proof of the exact Frozen Claim is complete, with no `sorry`, no additional axioms beyond Lean/Mathlib's standard ones, and no auxiliary hypotheses.
\item The narrative prose above (including the two Remarks, which are commentary rather than part of the theorem) has not been independently assessed; only the Lean statement and proof were kernel verified.
\item The generalization mentioned in the Remarks — that √a + √b is irrational whenever a, b are nonnegative rationals with √a irrational — was not formalized or verified; it is stated as informal commentary only.
\item The proof relies on Mathlib's `irrational\_sqrt\_two`, `Real.sq\_sqrt`, and the standard ring/field structure of ℚ and ℝ; these are taken as given from the pinned Mathlib revision rather than reproved here.
\end{itemize}

\section*{Exact Lean statement}
\begin{Verbatim}[fontsize=\small]
theorem irrational_sqrt_two_add_sqrt_three : Irrational (Real.sqrt 2 + Real.sqrt 3)
\end{Verbatim}

\section*{Verification appendix}
\textbf{Axioms used:} propext, Classical.choice, Quot.sound

\begin{Verbatim}[fontsize=\small]
Frozen Claim SHA-256: 371232aaf211dc1fa654c850779771e449297d01f99cc286d6fd1dd320cb1dfc
Lean source SHA-256: 056705a85205dc64acf5084aedad9ceac0fb9f0253a5eaca534d02230eaa473a
Verification SHA-256: 3083bfbb4a5dc6e6d043571a4a9909cae8ce115b8805662d16d63dcede1b143a
Prompt set SHA-256: c418edba223323cd3bbb730ae110d1270be3d23e93994f25a6ddcf2c7f9436a6
Tectonic bundle SHA-256: 6ffe055852f8faf66c0acbe1a7fb27f87b869a90bad1204f3bf4d9683f597c7c
\end{Verbatim}

\begin{samepage}
\textbf{Run and tool identities}
\begin{itemize}
\item Run ID: 8ccb35a8-74bd-4af0-9015-71cf0b1dedca
\item Model: claude-opus-5
\item Backend: claude
\item Runtime SDK: 0.2.127
\item Lean: 4.33.1
\item Mathlib: 0df444a360eaa60ab8c11dca51a86af692955474
\item Tectonic: 0.16.9
\item Tectonic bundle: https://data1.fullyjustified.net/tlextras-2022.0r0.tar
\end{itemize}
\end{samepage}

\vfill
\small
Lean kernel verification, faithfulness (approved by the user and read back by an
independent model before proof search), heuristic informal review, and document
compilation are independent statuses. This trusted-development MVP
does not sandbox generated Lean or TeX.

\end{document}
